\begin{abstract}
The rapid proliferation of large language models (LLMs) complicates deployment decisions, as performance varies significantly across models, prompts, and cost regimes. Static LLM routers trained on offline preference data often degrade when deployment distributions drift or prompt-specific model preferences differ from training assumptions. In our evaluation, 14.1\% of prompts favor a cheaper model over a more expensive frontier model, causing static threshold-based routers to behave non-monotonically and overspend. To address this, we introduce banditGPT, a cost-aware contextual bandit framework for dynamic LLM routing. banditGPT employs online learning to adapt to deployment distributions without requiring supervised labels. By integrating a Corralling meta-learner to balance offline warmup priors with tabula rasa exploration, and a Hybrid LinUCB mechanism to pool observations across data-driven model families, the system rapidly calibrates to new distributions. Experimental results demonstrate that banditGPT closes 70.0\% of the optimality gap to an oracle router within 400 online interactions. Across portfolios spanning six providers and a 600$\times$ cost range, banditGPT achieves up to an 89\% cost reduction with less than 6\% quality degradation. Ablation studies indicate that continuous family-level sharing provides an effective abstraction for cross-model knowledge transfer, while one-shot semantic transfer adds no significant benefit. Code, warmup priors, and experimental logs are open-sourced to facilitate further research in adaptive LLM routing.
\end{abstract}
